\documentclass[10pt,fleqn]{article}
\usepackage{hyperref}
\usepackage{graphicx}

\setlength{\topmargin}{-.75in}
\addtolength{\textheight}{2.00in}
\setlength{\oddsidemargin}{.00in}
\addtolength{\textwidth}{.75in}

\nofiles

\pagestyle{empty}

\setlength{\parindent}{0in}

\input{/cygdrive/m/tex/commands/commands}

\begin{document}
%%%%%%%%%%%%%%%%%%%%%%%%%%%%%%%%%%%%%%%%%%%%%%%%%%%%%%%%%%%%%%%%%%%%%%%%%%%%%%%%
%%%%%%%%%%%%%%%%%%%%%%%%%%%%%%%%%%%%%%%%%%%%%%%%%%%%%%%%%%%%%%%%%%%%%%%%%%%%%%%%
\vskip0.1in\hrule\vskip0.1in \noindent
{\bf{\Large Math 4610 Fundamentals of Numerical Analysis Tasksheet 6}}
\vskip0.1in\hrule\vskip0.1in \noindent
The problems for the Tasksheet 6 are included below. The deadline for turning
in your work on these problems will be posted on the repository. In addition,
you will turn in your work through the math4610 repository. A directory will be
constructed that will be used as a place to store your work.
\vskip0.1in\hrule\vskip0.1in \noindent
{\bf{\large Tasks}}
\vskip0.1in\hrule\vskip0.1in \noindent
\begin{trivlist}
  \item[\bf Task 1:]
    In this problem consider the function
    \[
      f(x) = e^{-x^2}\ sin(4\ x^2-1.0) + 0.051
    \]
    Try all four methods for finding roots that you have developed. Start by
    computing the closest root to zero. Hint: Start with Bisection on a small
    interval around \(x^*=0.0\) to bracket the root and then apply the other
    methods to verify they work. Write a code that is linked to your shared
    library, jar-file, or python module folder.
\vskip0.1in\hrule\vskip0.1in \noindent
  \item[\bf Task 2:] Apply your Newton method to the function in Task 1 by
    starting \(x_0=-5.0\) and \(x_0=6\). If we want the root closest to zero, as
    in Task1, compare the results obtained to the root we found in the previous
    Task. What went right or wrong?
\vskip0.1in\hrule\vskip0.1in \noindent
  \item[\bf Task 3:] Come up with an algorithm that starts by applying something
    like a hybrid method starting with an initial interval, \([-5,6]\), that
    will find the closest root to zero. Hint: you will need to modify the logic
    a little in the Bisection part to avoid finding different zeros than the one
    found in Task 1.
\vskip0.1in\hrule\vskip0.1in \noindent
  \item[\bf Task 4:] Repeat Task 3 for the hybrid method using the secant
    method. Use the results to verify the results in Task 3.
\vskip0.1in\hrule\vskip0.1in \noindent
  \item[\bf Task 5:] Let's change the problem a bit. Use your hybrid code to
    find as many roots as possible in the interval, \([-5,6]\)/. Note that you
    will need to write a bit of logic into your code that breaks the interval
    into intervals acceptable for use of the Bisection method. Hint: Write code
    that will produce subintervals where \(f(x_k)\ f(x_{k+1})<0\). You may need
    to refine the points which means dividing the initial interval into smaller
    and smaller intervals.
\vskip0.1in\hrule\vskip0.1in \noindent
  \item[\bf Task 6:] Search the internet for sites that discuss location of
    multiple roots in one dimension. Discuss the methods used to find additional
    toots. Write a brief summary of what you find including the pros and cons of
    the methods. Your write up should be a brief paragraph (3 or 4 sentences)
    that describe your findings. Include links to the sites you cite.
\end{trivlist}
\vskip0.1in\hrule\vskip0.1in \noindent
  \href{../../tasksheet_05/html/tasksheet_05.html}{Previous}
  \href{../../toc/md/tasksheet_toc.md}{Table of Contents} |
  \href{../../tasksheet_07/html/tasksheet_07.html}{Next}
\vskip0.1in\hrule\vskip0.1in \noindent
%%%%%%%%%%%%%%%%%%%%%%%%%%%%%%%%%%%%%%%%%%%%%%%%%%%%%%%%%%%%%%%%%%%%%%%%%%%%%%%%
%%%%%%%%%%%%%%%%%%%%%%%%%%%%%%%%%%%%%%%%%%%%%%%%%%%%%%%%%%%%%%%%%%%%%%%%%%%%%%%%
\end{document}
